\section{The intrinsic monomial ideal for a disconnected SNC boundary}
\label{sec:hironaka-snc-supported-monomial-ideal}

The members of a finite SNC list need not be connected.  Their divisor
exponents are locally constant on their smooth carriers, and can differ
between connected components.  The intrinsic monomial ideal uses these
actual stalk exponents and remains globally defined.  The following
argument identifies it with the actual ideal of every effective divisor
supported on the boundary, without assuming any global multiplicity for
a disconnected member.  It implements the monomial/nonmonomial
decomposition of W{\l}odarczyk's Step~2 at the level of genuine ideals.

\begin{lemma}[The intrinsic monomial ideal is invertible]
\label{hir:intrinsic-monomial-invertible}
\uses{hir:monomial-part-ambient}
\lean{Hironaka.HironakaSNCSupportedMonomialIdeal.isInvertible_monomialPart}
\leanok
Let $n>0$, let $E$ be a finite SNC list, and let $\mathfrak a$ be a
coherent ideal with a nonzero stalk at every point.  Its intrinsic
monomial ideal $\mathcal M(\mathfrak a)$ is invertible.
\end{lemma}
\begin{proof}
The existing genuine monomial generator theorem supplies, near each
$x$, an open set $U$ and a holomorphic function $m$ generating every
stalk of $\mathcal M(\mathfrak a)$ on $U$.  At each $y\in U$, points
where $m$ is nonzero occur arbitrarily close to $y$.
If $V\subseteq U$ is a nonempty open set, choose $y\in V$; the latter
property supplies $z\in V$ with $m(z)\ne0$.
Thus $m$ vanishes identically on no nonempty open subset of $U$.
These are exactly the local generator and nonvanishing properties of
an invertible holomorphic ideal.
\end{proof}

\begin{lemma}[The intrinsic nonmonomial ideal is the actual colon]
\label{hir:intrinsic-nonmonomial-colon}
\uses{hir:monomial-part-ambient}
\lean{Hironaka.HironakaSNCSupportedMonomialIdeal.nonmonomialPart_eq_colon}
\leanok
For every germ ideal $\mathfrak a$ and finite boundary list $E$,
\[
 \mathcal N(\mathfrak a)=
 (\mathfrak a:\mathcal M(\mathfrak a)).
\]
\end{lemma}
\begin{proof}
Both sides consist, at every point, of holomorphic germs $f$ such that
$fm$ belongs to $\mathfrak a$ for every germ
$m\in\mathcal M(\mathfrak a)$.  The product $fm=mf$ identifies the two
membership conditions, and hence the complete germ ideals.
\end{proof}

\begin{lemma}[The nonmonomial part of an effective divisor is invertible]
\label{hir:intrinsic-divisor-nonmonomial-invertible}
\uses{hir:intrinsic-monomial-invertible,hir:intrinsic-nonmonomial-colon,hir:monomial-part-ambient,hir:divisor-ideal-nonzero,ana:actual-invertible-ideal-quotient}
\lean{Hironaka.HironakaSNCSupportedMonomialIdeal.isInvertible_nonmonomialPart}
\leanok
If $n>0$ and $E$ is a finite SNC list, the intrinsic nonmonomial ideal
of an effective divisor $D$ is an actual invertible ideal.
\end{lemma}
\begin{proof}
The ideal $I_D$ is invertible, hence coherent with nonzero stalks.
The first lemma makes $\mathcal M(I_D)$ invertible, and the existing
divisor-exponent and Hadamard calculus gives
$I_D\subseteq\mathcal M(I_D)$ at all original points.
The actual colon of two nested invertible ideals is invertible: locally
their generators have a holomorphic quotient, and cancellation against
the nonzero divisor generator proves that this quotient generates the
entire colon stalk.  Apply that theorem and the preceding colon identity.
\end{proof}

\begin{lemma}[An SNC boundary has no point in dimension zero]
\label{hir:snc-dimension-zero-empty}
\uses{hir:snc}
\lean{Hironaka.HironakaSNCSupportedMonomialIdeal.notMem_member_of_dim_zero,Hironaka.HironakaSNCSupportedMonomialIdeal.monomialPart_eq_unitIdeal_of_dim_zero}
\leanok
On a complex zero-dimensional manifold, every member of a finite SNC
list is empty, and the intrinsic monomial ideal of any ideal is the
holomorphic unit ideal.
\end{lemma}
\begin{proof}
If $x$ belonged to a member, an adapted SNC chart would cut that member
out by a set of exactly one coordinate.  There are no coordinates in
dimension zero, so such a coordinate set cannot exist.
Thus there is no active member at any point.  The definition of the
intrinsic monomial ideal consequently imposes only holomorphicity on
its germs, which is the unit ideal.
\end{proof}

\begin{theorem}[The genuine nonmonomial residual is the unit ideal]
\label{hir:intrinsic-supported-nonmonomial-unit}
\uses{hir:intrinsic-divisor-nonmonomial-invertible,hir:divisor-supported-local-component,hir:monomial-part-ambient,hir:snc-dimension-zero-empty}
\dcref{W{\l}odarczyk 2005, Step 2}
\source{wlodarczyk05-hironaka}
\lean{Hironaka.HironakaSNCSupportedMonomialIdeal.nonmonomialPart_eq_unitIdeal_of_support_subset}
\leanok
Let $E$ be a finite SNC list on a complex manifold of any dimension,
and let $D$ be an effective divisor with
$\operatorname{Supp}D\subseteq\bigcup_{W\in E}W$.
Then its complete intrinsic nonmonomial ideal is the unit ideal:
\[
 \mathcal N(I_D)=\mathcal O_M.
\]
No connectedness of the members is required.
\end{theorem}
\begin{proof}
First suppose $n>0$.  The nonmonomial ideal is invertible by the previous
lemma, so it defines an actual effective divisor $R$.
The genuine factorization
$I_D=\mathcal N(I_D)\mathcal M(I_D)$ implies
$\operatorname{Supp}R\subseteq\operatorname{Supp}D$:
if all germs of its first factor vanish at a point, all germs of the
product vanish there as well.
If $x\in\operatorname{Supp}R$, the actual local component producer
gives $W\in E$ through $x$ with $W$ locally contained in
$\operatorname{Supp}R$.
The existing maximal-exponent extraction theorem supplies a germ
$f\in\mathcal N(I_D)_x$ which does not belong to $(I_W)_x$.
Coherence propagates membership of this same germ $f$ to nearby points.
At every such point in $W$, local support containment therefore gives
$f=0$.  Since $f$ is holomorphic, this says precisely
$f\in(I_W)_x$, a contradiction.
Thus $R$ has empty support, and a coherent ideal with empty zero locus
is the unit ideal, as proved by dividing by a nonvanishing stalk member.

When $n=0$, the preceding lemma shows that the boundary union is empty,
so the actual support of $D$ is empty and $I_D=\mathcal O_M$.
Every holomorphic germ then belongs to the nonmonomial ideal: its
product with any holomorphic monomial-ideal germ still belongs to
$\mathcal O_M$.  Conversely nonmonomial-ideal germs are holomorphic by
definition.  This proves the asserted equality in dimension zero too.
\end{proof}

\begin{theorem}[The actual supported ideal equals its intrinsic monomial ideal]
\label{hir:intrinsic-supported-ideal-monomial}
\uses{hir:intrinsic-supported-nonmonomial-unit,hir:monomial-part-ambient,hir:snc-dimension-zero-empty}
\lean{Hironaka.HironakaSNCSupportedMonomialIdeal.ideal_eq_monomialPart_of_support_subset}
\leanok
For every finite SNC list $E$ and every effective divisor $D$ with
$\operatorname{Supp}D\subseteq\bigcup_{W\in E}W$, there is equality of
the complete original germ ideals
\[
 I_D=\mathcal M(I_D).
\]
\end{theorem}
\begin{proof}
For positive ambient dimension, substitute the proved unit residual
into the genuine factorization
$I_D=\mathcal N(I_D)\mathcal M(I_D)$.
The monomial ideal is coherent, so multiplication by the holomorphic
unit ideal leaves it unchanged.  In dimension zero, the boundary is
empty, hence the supported ideal and its monomial ideal are both the
unit ideal.  This proves the equality at all original stalks in every
dimension.  The statement permits disconnected members: their actual
locally constant exponents are retained in $\mathcal M(I_D)$, rather
than replaced by unjustified globally fixed integers.
\end{proof}
